
\documentclass{article}
\usepackage{colortbl}
\usepackage{makecell}
\usepackage{multirow}
\usepackage{supertabular}

\begin{document}

\newcounter{utterance}

\centering \large Interaction Transcript for game `eqbench', experiment `eq\_bench', episode 49 with reflexion{-}Qwen3.5{-}9B.
\vspace{24pt}

{ \footnotesize  \setcounter{utterance}{1}
\setlength{\tabcolsep}{0pt}
\begin{supertabular}{c@{$\;$}|p{.15\linewidth}@{}p{.15\linewidth}p{.15\linewidth}p{.15\linewidth}p{.15\linewidth}p{.15\linewidth}}
    \# & \multicolumn{2}{c}{Player} && \multicolumn{2}{c}{Game Master} \\
    \hline

    \theutterance \stepcounter{utterance}  
    & & & \multicolumn{4}{p{0.6\linewidth}}{
        \cellcolor[rgb]{0.9,0.9,0.9}{
            \makecell[{{p{\linewidth}}}]{
                \texttt{\tiny{[P1$\langle$GM]}}
                \texttt{You have attempted this same task before and did not succeed. These are the plans you wrote after those attempts.} \\
\texttt{Trial \#1: New plan: I incorrectly prioritized Jen's aggressive tone as the sole indicator of her emotion, leading me to overestimate "Defiant" and underestimate "Frustrated." I failed to recognize that Dr. Hall's final statement ("It's a classroom") directly addresses Jen's refusal to accept his authority, which is the definition of being reprimanded. Next time, I will analyze the immediate context of the final line: since Jen was told she is in a classroom, not a courtroom, and she persisted in her debate stance, the primary emotional response to this correction is "Reprimanded." I will assign the highest score to Reprimanded (10) and lower scores to Defiant (5) and Frustrated (4), while keeping Dismissive at 0, as her engagement proves she is not dismissive.} \\
\\ 
\texttt{Trial \#2: New plan: I failed to recognize that Jen's aggressive retorts were a direct response to Dr. Hall's accusation that she was "overcompensating." Her tone was not merely defiance; it was a defensive reaction to being called out for her behavior. Dr. Hall's final line, "It's a classroom," serves as a formal correction of her stance, effectively reprimanding her for treating the setting as a debate arena. Therefore, the primary emotion must be "Reprimanded" (10), as she has been explicitly corrected by an authority figure. "Defiant" should be lower (5) because her defiance is reactive, not the primary state after the correction. "Frustrated" (4) is secondary to the conflict, and "Dismissive" (0) is incorrect given her active engagement. I will adjust the scores to reflect that the reprimand is the dominant outcome of the interaction.} \\
\\ 
\texttt{Reprimanded: 10} \\
\texttt{Dismissive: 0} \\
\texttt{Frustrated: 4} \\
\texttt{Defiant: 5} \\
\\ 
\texttt{Trial \#3: New plan: I failed to distinguish between the *cause* of the emotion (Jen's defiance) and the *result* of the interaction (being reprimanded). While Jen is defiant throughout the dialogue, the specific question asks for her feeling "at the end," which is defined by Dr. Hall's final authoritative correction ("It's a classroom"). This action transforms the dynamic from a mutual debate into a disciplinary moment. Therefore, "Reprimanded" must be the dominant emotion (10), as it is the direct consequence of the final line. "Defiant" should be reduced to a secondary trait (5) representing her underlying attitude, and "Frustrated" lowered to a minor reaction (2) to the interruption, while "Dismissive" remains 0 due to her active engagement.} \\
\\ 
\texttt{Reprimanded: 10} \\
\texttt{Dismissive: 0} \\
\texttt{Frustrated: 2} \\
\texttt{Defiant: 5} \\
\\ 
\\ 
\\ 
\texttt{Your task is to predict the likely emotional responses of a character in this dialogue:} \\
\\ 
\texttt{Dr. Hall: "Ms. Vasquez, I see you're gracing us with your presence today. I trust you've done the reading?"} \\
\texttt{Jen: "Oh, I wouldn't miss it for the world, Dr. Hall. And yes, I've done the reading. Twice."} \\
\texttt{Dr. Hall: "Twice? My, my. Overachieving or overcompensating, I wonder?"} \\
\texttt{Jen: "Funny you should mention overcompensation, Dr. Hall."} \\
\texttt{Dr. Hall: "Careful, Ms. Vasquez. Your wit is as sharp as a butter knife."} \\
\texttt{Jen: "And your condescension is as subtle as a sledgehammer, Dr. Hall."} \\
\texttt{Dr. Hall: "I see. You're here to debate, not to learn."} \\
\texttt{Jen: "I'm here to do both. Can't say the same for everyone, though."} \\
\texttt{Dr. Hall: "Enough. This is not a courtroom, Ms. Vasquez. It's a classroom."} \\
\texttt{{[}End dialogue{]}} \\
\\ 
\texttt{At the end of this dialogue, Jen would feel...} \\
\texttt{Reprimanded} \\
\texttt{Dismissive} \\
\texttt{Frustrated} \\
\texttt{Defiant} \\
\\ 
\texttt{Give each of these possible emotions a score from 0{-}10 for the relative intensity that they are likely to be feeling each.} \\
\\ 
\texttt{You must output in the following format, including headings (of course, you should give your own scores), with no additional commentary:} \\
\\ 
\texttt{Reprimanded: <score>} \\
\texttt{Dismissive: <score>} \\
\texttt{Frustrated: <score>} \\
\texttt{Defiant: <score>} \\
\\ 
\\ 
\texttt{{[}End of answer{]}} \\
\\ 
\texttt{Remember: zero is a valid score, meaning they are likely not feeling that emotion. You must score at least one emotion > 0.} \\
\\ 
\texttt{Your answer:} \\
            }
        }
    }
     \\ \\

    \theutterance \stepcounter{utterance}  
    & \multicolumn{4}{p{0.6\linewidth}}{
        \cellcolor[rgb]{0.561,0.737,0.561}{
            \makecell[{{p{\linewidth}}}]{
                \texttt{\tiny{[P1$\rangle$GM]}}
                \texttt{Reprimanded: 10} \\
\texttt{Dismissive: 0} \\
\texttt{Frustrated: 3} \\
\texttt{Defiant: 6} \\
            }
        }
    }
    & & \\ \\

    \theutterance \stepcounter{utterance}  
    & & & \multicolumn{2}{p{0.3\linewidth}}{
        \cellcolor[rgb]{0.9,0.9,0.9}{
            \makecell[{{p{\linewidth}}}]{
                \texttt{\tiny{[GM$|$GM]}}
                \texttt{Reprimanded: 10} \\
\texttt{Dismissive: 0} \\
\texttt{Frustrated: 3} \\
\texttt{Defiant: 6} \\
            }
        }
    }
    & & \\ \\

    \theutterance \stepcounter{utterance}  
    & & & \multicolumn{2}{p{0.3\linewidth}}{
        \cellcolor[rgb]{0.9,0.9,0.9}{
            \makecell[{{p{\linewidth}}}]{
                \texttt{\tiny{[GM$|$GM]}}
                \texttt{Reprimanded: 6} \\
\texttt{Dismissive: 7} \\
\texttt{Frustrated: 8} \\
\texttt{Defiant: 9} \\
            }
        }
    }
    & & \\ \\

    \theutterance \stepcounter{utterance}  
    & & & \multicolumn{2}{p{0.3\linewidth}}{
        \cellcolor[rgb]{0.9,0.9,0.9}{
            \makecell[{{p{\linewidth}}}]{
                \texttt{\tiny{[GM$|$GM]}}
                \texttt{game\_result = LOSE} \\
            }
        }
    }
    & & \\ \\

    \theutterance \stepcounter{utterance}  
    & & \multicolumn{4}{p{0.6\linewidth}}{
        \cellcolor[rgb]{0.95,0.95,0.95}{
            \makecell[{{p{\linewidth}}}]{
                \texttt{\tiny{[P1$\leftrightarrow$P]}}
                \texttt{Post{-}episode Reflexion (not sent to GM):} \\
\texttt{New plan: I consistently failed to lower the score for "Frustrated" enough, likely because I focused too heavily on Jen's active verbal sparring (defiance) and the authority figure's correction (reprimand), while underestimating the cumulative effect of Dr. Hall's condescending remarks ("sharp as a butter knife," "not a courtroom") on Jen's emotional state. The dialogue shows a clear escalation where Jen's wit is met with sarcasm and dismissal, which naturally breeds frustration before the final reprimand. To correct this, I must recognize that while "Reprimanded" is the dominant outcome of the final line, "Frustrated" is the primary driver of her behavior throughout the exchange and should be scored higher than "Defiant." I will assign "Reprimanded" the highest score (10) as the immediate reaction to the final line, but increase "Frustrated" to 8 to reflect the sustained irritation from the tone of the entire conversation, and lower "Defiant" to 4 as it is a coping mechanism rather than the primary feeling. "Dismissive" remains 0.} \\
\\ 
\texttt{Reprimanded: 10} \\
\texttt{Dismissive: 0} \\
\texttt{Frustrated: 8} \\
\texttt{Defiant: 4} \\
            }
        }
    }
    & \\ \\

\end{supertabular}
}

\end{document}
